\section{Threshold MPC vs.\ On-Chain Multi-Signature}

\label{sec:mpc-vs-multisig}
%% =====================================================================

The choice between threshold multi-party-computation signing and
on-chain multi-signature constructions is the most consequential
single design decision in this architecture. This section makes that
decision explicit and defends it.

\subsection{The Two Constructions}

\paragraph{On-chain multi-signature.} The wallet exposes a public
script (Bitcoin: P2WSH or Taproot script-path; Ethereum: a smart
contract such as Gnosis Safe). The script encodes the threshold
$t$-of-$n$ and the public keys of all $n$ signers. Spending requires
$t$ valid signatures, each from a distinct on-chain key, presented in
the witness or contract call.

\paragraph{Threshold MPC.} The wallet exposes a single public key.
The corresponding private key has never existed in one place: it was
generated as $n$ shares via a distributed key-generation protocol,
and is used by an interactive multi-party signing protocol that
produces a single canonical chain signature. The chain sees one
signature; the threshold structure is invisible.

\subsection{Comparison}

\begin{table}[h]
\centering
\small
\begin{tabular}{p{4.5cm}p{5cm}p{5cm}}
\toprule
\textbf{Property} & \textbf{On-chain multisig} & \textbf{Threshold MPC} \\
\midrule
Privacy & Threshold and signers are visible on chain &
Wallet looks like a single-key wallet \\
Chain fees & 2--5x base txn (Bitcoin P2WSH) or 30--80k gas extra
            (Gnosis Safe) & Single-sig fees \\
Cross-chain consistency & Each chain has different multisig semantics;
audit and policy must be reimplemented per chain &
One signing pipeline; same code for BTC and ETH \\
Throughput & Limited by multisig contract gas or witness size &
Limited by MPC round-trip latency only \\
Forward security & Compromise reveals which signer signed &
Aggregate signature carries no per-share attribution \\
Policy disclosure & The set of signers and the threshold are public &
Policy lives off-chain \\
On-chain code surface & Smart-contract wallet is itself attack
surface (cf.\ Parity Multisig 2017 freeze) &
No smart-contract wallet \\
Quantum hedge integration & Difficult; contract-bound to ECDSA &
Can be paired with M-Chain ML-DSA audit log (LP-006 \cite{lp006}) \\
\bottomrule
\end{tabular}
\caption{Multisig vs.\ threshold MPC.}
\label{tab:mpc-vs-multisig}
\end{table}

\subsection{Why MPC for an Eight-Billion-Dollar Fund}

Three properties dominate the decision.

\paragraph{Privacy of policy.} An on-chain $3$-of-$5$ Bitcoin wallet
publishes, in every spend, the five public keys and the threshold.
An adversary planning a targeted attack now knows exactly which five
signers to compromise and exactly the threshold to clear. With
threshold MPC the same fund publishes one address and one signature;
the adversary observes nothing about the custody policy. For a
public-figure, sovereign, or large-pension custodian this is
materially important.

\paragraph{Cross-chain uniformity.} Bitcoin multisig (P2WSH or
Taproot script-path), Ethereum multisig (Gnosis Safe), and the
multisig idioms of the Lux X-Chain and C-Chain do not share semantics.
Each chain's multisig has its own fee profile, its own upgradability
risk, and its own auditing surface. Threshold MPC over secp256k1
produces a single canonical signature acceptable to every secp256k1
chain; one signing pipeline in \texttt{luxfi/mpc} signs Bitcoin,
Ethereum, all EVM chains, and the Lux X-Chain identically.

\paragraph{Future-state quantum-resistant migration.} When the
post-quantum migration window opens, a fund holding ECDSA-signed
addresses must roll its custody. The migration target on Lux is the
M-Chain BLS+ML-DSA hybrid signing path (LP-006 \cite{lp006},
LP-020 \cite{lp020}). Threshold MPC's chain-invisibility means that
the migration is operationally a key rotation; on-chain multisig
contracts would need a contract-level migration with all the upgrade
risk that entails.

\subsection{The CGGMP21 Choice for ECDSA}

ECDSA threshold signing is harder than threshold Schnorr or
threshold BLS because ECDSA's nonce inversion is non-linear in the
secret. We adopt CGGMP21 \cite{cggmp21} for the following reasons.

\begin{enumerate}[nosep]
\item \textbf{UC-secure with identifiable abort.} A misbehaving
party is identified, not merely detected, which is a hard requirement
when shares are owned by independent legal entities.
\item \textbf{Non-interactive online phase.} Pre-signing computes
ahead of time; the live signing latency is dominated by network
round-trips rather than zero-knowledge proofs.
\item \textbf{Proactive refresh.} Shares can be refreshed without
changing the public key, which lets us rotate signing-side credentials
on the regular operational cadence (Section~\ref{sec:ops}) without
emitting a chain transaction.
\item \textbf{Implementations are public and audited.} The
\texttt{luxfi/mpc} implementation is the canonical Go codepath; it
is audited and has a published test-vector suite.
\end{enumerate}

\subsection{The FROST Path for Schnorr}

Where the chain accepts Schnorr signatures (Bitcoin Taproot
key-spend, Lux M-Chain BLS-aggregate path), FROST \cite{frost} is
preferred over CGGMP21:

\begin{itemize}[nosep]
\item Two-round signing with a one-round commit phase.
\item No range proofs and no MtA blowup.
\item Native Schnorr (BIP-340) so the on-chain signature is
indistinguishable from a single-key Taproot key-spend.
\end{itemize}

For the present architecture FROST is used on Bitcoin Taproot
key-paths and CGGMP21 is used on Ethereum and EVM chains. Both
codepaths share the same ceremony, attestation, and audit-log
surfaces; only the signing primitive differs.

\subsection{Why Not Both?}

A defensible alternative is to run on-chain multisig as a backup
recovery vehicle (e.g.\ a $3$-of-$5$ on-chain multisig at the cold
tier, signed off only when the MPC pipeline is unavailable). We
evaluated this and rejected it:

\begin{itemize}[nosep]
\item It doubles the audit surface and the ceremony surface.
\item The on-chain multisig leaks the cold-tier signing graph at
the moment it is published, even if it is never used.
\item The same disaster-recovery property is achieved by the
sealed-envelope shard scheme (Section~\ref{sec:dr}) without any
on-chain footprint.
\end{itemize}

\subsection{Conclusion}

Threshold MPC is selected for all three tiers. The custody policy
is invisible to the chain. The signing pipeline is one pipeline
across BTC, ETH, and Lux. The migration to a post-quantum signing
primitive is a key rotation rather than a contract migration.

%% =====================================================================
